\section{Ring maps of finite type and of finite presentation}
\label{section-finite-type}

\begin{definition}
\label{definition-finite-type}
Let $R \to S$ be a ring map.
\begin{enumerate}
\item We say $R \to S$ is of {\it finite type}, or that {\it $S$ is a finite
type $R$-algebra} if there exist an $n \in \mathbf{N}$ and an surjection
of $R$-algebras $R[x_1, \ldots, x_n] \to S$.
\item We say $R \to S$ is of {\it finite presentation} if there
exist integers $n, m \in \mathbf{N}$ and polynomials
$f_1, \ldots, f_m \in R[x_1, \ldots, x_n]$
and an isomorphism of $R$-algebras
$R[x_1, \ldots, x_n]/(f_1, \ldots, f_m) \cong S$.
\end{enumerate}
\end{definition}

\noindent
Informally, $R \to S$ is of finite presentation if and only if
$S$ is finitely generated as an $R$-algebra
and the ideal of relations among the generators is finitely generated.
A choice of a surjection $R[x_1, \ldots, x_n] \to S$ as in the definition
is sometimes called a {\it presentation} of $S$.

\begin{lemma}
\label{lemma-compose-finite-type}
The notions finite type and finite presentation have the following
permanence properties.
\begin{enumerate}
\item A composition of ring maps of finite type is of finite type.
\item A composition of ring maps of finite presentation is of finite
presentation.
\item Given $R \to S' \to S$ with $R \to S$ of finite type,
then $S' \to S$ is of finite type.
\item Given $R \to S' \to S$, with $R \to S$ of finite presentation,
and $R \to S'$ of finite type, then $S' \to S$ is of finite presentation.
\end{enumerate}
\end{lemma}

\begin{proof}
We only prove the last assertion.
Write $S = R[x_1, \ldots, x_n]/(f_1, \ldots, f_m)$
and $S' = R[y_1, \ldots, y_a]/I$. Say that the class
$\bar y_i$ of $y_i$ maps
to $h_i \bmod (f_1, \ldots, f_m)$ in $S$.
Then it is clear that
$S = S'[x_1, \ldots, x_n]/(f_1, \ldots, f_m,
h_1 - \bar y_1, \ldots, h_a - \bar y_a)$.
\end{proof}

\begin{lemma}
\label{lemma-finite-presentation-independent}
Let $R \to S$ be a ring map of finite presentation.
For any surjection $\alpha : R[x_1, \ldots, x_n] \to S$ the
kernel of $\alpha$ is a finitely generated ideal in $R[x_1, \ldots, x_n]$.
\end{lemma}

\begin{proof}
Write $S = R[y_1, \ldots, y_m]/(f_1, \ldots, f_k)$.
Choose $g_i \in R[y_1, \ldots, y_m]$ which are lifts
of $\alpha(x_i)$. Then we see that $S = R[x_i, y_j]/(f_l, x_i - g_i)$.
Choose $h_j \in R[x_1, \ldots, x_n]$ such that $\alpha(h_j)$
corresponds to $y_j \bmod (f_1, \ldots, f_k)$. Consider
the map $\psi : R[x_i, y_j] \to R[x_i]$, $x_i \mapsto x_i$,
$y_j \mapsto h_j$. Then the kernel of $\alpha$
is the image of $(f_l, x_i - g_i)$ under $\psi$ and we win.
\end{proof}

\begin{lemma}
\label{lemma-finitely-presented-over-subring}
Let $R \to S$ be a ring map.
Let $M$ be an $S$-module.
Assume $R \to S$ is of finite type and
$M$ is finitely presented as an $R$-module.
Then $M$ is finitely presented as an $S$-module.
\end{lemma}

\begin{proof}
This is similar to the proof of part (4) of
Lemma \ref{lemma-compose-finite-type}.
We may assume $S = R[x_1, \ldots, x_n]/J$.
Choose $y_1, \ldots, y_m \in M$ which generate $M$ as an $R$-module
and choose relations $\sum a_{ij} y_j = 0$, $i = 1, \ldots, t$ which
generate the kernel of $R^{\oplus m} \to M$. For any
$i = 1, \ldots, n$ and $j = 1, \ldots, m$ write
$$
x_i y_j = \sum a_{ijk} y_k
$$
for some $a_{ijk} \in R$. Consider the $S$-module $N$ generated by
$y_1, \ldots, y_m$ subject to the relations
$\sum a_{ij} y_j = 0$, $i = 1, \ldots, t$ and
$x_i y_j = \sum a_{ijk} y_k$, $i = 1, \ldots, n$ and $j = 1, \ldots, m$.
Then $N$ has a presentation
$$
S^{\oplus nm + t} \longrightarrow S^{\oplus m} \longrightarrow N
\longrightarrow 0
$$
By construction there is a surjective map $\varphi : N \to M$.
To finish the proof we show $\varphi$ is injective.
Suppose $z = \sum b_j y_j \in N$ for some $b_j \in S$.
We may think of $b_j$ as a polynomial in $x_1, \ldots, x_n$
with coefficients in $R$.
By applying the relations of the form $x_i y_j = \sum a_{ijk} y_k$
we can inductively lower the degree of the polynomials.
Hence we see that $z = \sum c_j y_j$ for some $c_j \in R$.
Hence if $\varphi(z) = 0$ then the vector $(c_1, \ldots, c_m)$
is an $R$-linear combination of the vectors $(a_{i1}, \ldots, a_{im})$
and we conclude that $z = 0$ as desired.
\end{proof}
